\section{Reproduction}\label{app:repro}

All paths are relative to the code release (\url{https://github.com/kjungmo/se-mppi-nav2}, Apache-2.0).

\noindent\textbf{Software.} ROS~2 Jazzy, Nav2 and Gazebo are installed through RoboStack (conda channels \code{robostack-jazzy} and \code{conda-forge}) by \code{scripts/setup_ros2_env.sh}, which also installs \code{osqp-eigen} and \code{eigen} for the CBF-QP and a compiler toolchain, so the build does not depend on the host OS version.

\noindent\textbf{Benchmark regeneration.} Every benchmark number is generated from the committed per-trial file by committed scripts. The committed data use $N=50$:
{\small
\begin{verbatim}
export MAMBA_ROOT_PREFIX=$HOME/micromamba
micromamba run -n ros2 python3 -m experiments.benchmark2d.runner \
    --families utrap clutter dynamic narrowdyn -n 50 --workers 6 --max-steps 400
micromamba run -n ros2 python3 -m experiments.benchmark2d.report
\end{verbatim}
}
\noindent The figures and tables of this paper are regenerated from those outputs, and the mechanism runs of \cref{tab:mechanism} are re-executed, by
{\small
\begin{verbatim}
micromamba run -n ros2 python3 docs/papers/arxiv_ref/scripts/make_figures.py
python3 docs/papers/arxiv_ref/scripts/make_tables.py
python3 scripts/check_paper_numbers.py --arxiv
\end{verbatim}
}

\noindent\textbf{Scenarios and seeds.} Seeds $0$--$49$ per family; the generators in \code{experiments/benchmark2d/scenarios.py} seed \code{numpy} generators with a per-family offset plus the trial seed ($10^{6}$ U-trap, $2\times10^{6}$ clutter, $3\times10^{6}$ dynamic, $4\times10^{6}$ narrow-dynamic), and the MPPI noise stream uses the same trial seed (\code{experiments/benchmark2d/runner.py}), so every configuration sees identical worlds and rollout noise. The six toggle sets are in \code{experiments/benchmark2d/configs.py}; the overlays for the deferred full-scale protocol are in \code{experiments/configs/ablations.yaml}.

\noindent\textbf{Data.} Raw per-trial data: \code{experiments/results_2d/trials.csv} ($1{,}200$ rows). Aggregates and statistics: \code{summary.csv}, \code{stats.json}, \code{tables.md} in the same directory. Live pilot telemetry and logs: \code{experiments/results_pilot/}. The claim-to-artifact map for this paper is \code{docs/papers/arxiv_ref/NUMBERS.md}.

\section{Parameters}\label{app:params}

\begin{table}[!htbp]
\centering
\caption{\textbf{Parameters.} Released-controller defaults come from the package's parameter file and parameter getters; 2D-mirror values from the prototype primitives and the benchmark rollout (\cref{app:repro,app:mirror}).}
\label{tab:params}
\footnotesize
\setlength{\tabcolsep}{4pt}
\begin{tabular}{llcc}
\toprule
Symbol & Meaning & Released controller & 2D mirror \\
\midrule
$T,\ \Delta t$ & MPPI horizon steps, step (s) & 56, 0.05 & 25, 0.1 \\
$K$ & rollouts per cycle & 1000 & 600 \\
$\lambda$ & MPPI temperature & 0.3 & 0.3 \\
$\Sigma^{1/2}$ & sampling std.\ of $(v,\omega)$ & (0.2, 0.4) & (0.25, 0.6) \\
$\mathcal{U}$ & $v\in[v_{\min},v_{\max}]$, $|\omega|\le\omega_{\max}$ & $[-0.35, 0.5]$, 1.9 & $[-0.35, 0.5]$, 1.9 \\
$W$ & entrapment stall window (cycles) & 30 & 15 \\
$\epsilon_{\mathrm{goal}}$ & entrapment suppression radius (m) & 0.5 & -- \\
$d_0,\ \eta$ & APF influence distance (m), gain & 0.6, 1.0 & 0.8, 0.3 \\
$N_r,\ r_{\max},\ r_{\min}$ & gap rays, ray range (m), min.\ clearance (m) & 36, 2.0, 0.6 & 72, 3.0, 2.4 \\
$w_g$ & gap-attraction weight & 4.0 & subgoal (\cref{app:mirror}) \\
$L$ & look-ahead distance (m) & 0.2 & 0.2 \\
$r,\ m$ & robot radius (m), safety margin (m) & 0.225 (logged), 0.05 & 0.22, 0.05 \\
$\rho$ & QP slack weight & $10^{3}$ & $10^{3}$ \\
$\epsilon_\delta$ & slack threshold for braking & $10^{-3}$ & $10^{-3}$ \\
$N_{\mathrm{obs}}$ & obstacles kept in the QP & 20 & all dynamic \\
$\alpha_{\mathrm{base}},\ \alpha_{\mathrm{escape}}$ & CBF class-$\mathcal{K}$ gains & 2.0, 6.0 & 2.0, 6.0 \\
$\tau$ & TTC override threshold (s) & 1.5 & 1.5 \\
$v_{\mathrm{dyn}},\ R_{\max}$ & CBF admission: min.\ speed (m/s), max.\ radius (m) & 0.1, 1.0 & $\|\mathbf{v}_o\|>0$ \\
\bottomrule
\end{tabular}
\end{table}


\section{Differences Between the 2D Mirror and the Nav2 Controller}\label{app:mirror}

The 2D mirror (\code{experiments/prototype/se_mppi_proto.py}, driven by \code{experiments/benchmark2d/rollout.py}) reuses the controller's formulas but differs in the following documented ways.
\begin{itemize}
    \item \emph{Progress signal.} Without a global path, entrapment is declared when the best distance-to-goal has not improved by $10^{-3}$\,m for $W=15$ cycles, instead of the stall of $k^{\max}_t$ in \eqref{eq:detect}.
    \item \emph{Gap use.} The selected gap bearing places a temporary subgoal $2.5$\,m along the opening, replacing the goal in the MPPI cost while entrapped, in lieu of Nav2's global routing and the critic's $C_{\mathrm{gap}}$ term; the gap choice carries a hysteresis term toward the previous bearing so that symmetric openings do not flip every cycle.
    \item \emph{Geometry.} Obstacles, including walls built from chains of discs, are circles; APF distances and ray lengths are computed analytically against them rather than on a costmap distance field.
    \item \emph{Obstacle knowledge.} The CBF and TTC receive the true positions and velocities of the moving obstacles, so the tracker of \cref{sec:track} is not exercised.
    \item \emph{Costs.} The MPPI cost is a goal-distance, path-progress and obstacle-proximity cost rather than Nav2's critic stack; parameters differ as listed in \cref{tab:params}.
    \item \emph{Mechanism scenes.} \Cref{tab:mechanism} comes from \code{run_validation.py}, whose CBF sees all obstacles, including the static wall discs; the benchmark mirror restricts the CBF to movers, as the Nav2 controller does.
\end{itemize}

\section{Extended Results}\label{app:results}

\begin{table}[!htbp]
\centering
\caption{\textbf{Full per-family results} of the Python 2D-mirror benchmark ($N=50$ per cell). Rates in \%; time-to-goal and path length are means over successful trials only (\,$\pm$ 95\% CI half-width); minimum clearance is over all trials. Source: \texttt{experiments/results\_2d/summary.csv}.}
\label{tab:full}
\footnotesize
\setlength{\tabcolsep}{4pt}
\begin{tabular}{llcccccc}
\toprule
Family & Config & Succ. & Coll. & Timeout & Time-to-goal (s) & Path (m) & Min.\ clear.\ (m) \\
\midrule
\multirow{6}{*}{U-trap} & A & 0 & 0 & 100 & -- & -- & 0.33\stdv{0.00} \\
 & C & 90 & 0 & 10 & 31.34\stdv{1.36} & 8.12\stdv{0.33} & 0.32\stdv{0.00} \\
 & D & 0 & 0 & 100 & -- & -- & 0.33\stdv{0.00} \\
 & E & 88 & 0 & 12 & 31.20\stdv{1.36} & 8.09\stdv{0.32} & 0.32\stdv{0.00} \\
 & F & 88 & 0 & 12 & 31.20\stdv{1.36} & 8.09\stdv{0.32} & 0.32\stdv{0.00} \\
 & F$^{-}$ & 0 & 0 & 100 & -- & -- & 0.33\stdv{0.00} \\
\midrule
\multirow{6}{*}{Clutter} & A & 52 & 0 & 48 & 15.23\stdv{0.49} & 5.81\stdv{0.17} & 0.36\stdv{0.02} \\
 & C & 62 & 0 & 38 & 17.67\stdv{2.11} & 6.14\stdv{0.32} & 0.35\stdv{0.02} \\
 & D & 52 & 0 & 48 & 15.23\stdv{0.49} & 5.81\stdv{0.17} & 0.36\stdv{0.02} \\
 & E & 62 & 0 & 38 & 17.67\stdv{2.11} & 6.14\stdv{0.32} & 0.35\stdv{0.02} \\
 & F & 62 & 0 & 38 & 17.67\stdv{2.11} & 6.14\stdv{0.32} & 0.35\stdv{0.02} \\
 & F$^{-}$ & 52 & 0 & 48 & 15.23\stdv{0.49} & 5.81\stdv{0.17} & 0.36\stdv{0.02} \\
\midrule
\multirow{6}{*}{Dynamic} & A & 80 & 20 & 0 & 22.48\stdv{1.44} & 7.99\stdv{0.55} & 0.12\stdv{0.03} \\
 & C & 76 & 20 & 4 & 21.87\stdv{1.36} & 7.75\stdv{0.53} & 0.13\stdv{0.03} \\
 & D & 82 & 16 & 2 & 22.74\stdv{1.52} & 8.02\stdv{0.65} & 0.13\stdv{0.03} \\
 & E & 78 & 18 & 4 & 21.77\stdv{1.25} & 7.56\stdv{0.47} & 0.13\stdv{0.03} \\
 & F & 78 & 18 & 4 & 21.78\stdv{1.25} & 7.57\stdv{0.47} & 0.13\stdv{0.03} \\
 & F$^{-}$ & 84 & 16 & 0 & 23.01\stdv{1.66} & 8.14\stdv{0.68} & 0.13\stdv{0.03} \\
\midrule
\multirow{6}{*}{Narrow-dyn.} & A & 0 & 18 & 82 & -- & -- & 0.22\stdv{0.04} \\
 & C & 78 & 18 & 4 & 30.52\stdv{1.30} & 8.33\stdv{0.44} & 0.21\stdv{0.04} \\
 & D & 0 & 32 & 68 & -- & -- & 0.19\stdv{0.04} \\
 & E & 64 & 32 & 4 & 30.55\stdv{1.50} & 8.24\stdv{0.55} & 0.18\stdv{0.04} \\
 & F & 62 & 32 & 6 & 30.36\stdv{1.51} & 8.11\stdv{0.50} & 0.18\stdv{0.04} \\
 & F$^{-}$ & 0 & 32 & 68 & -- & -- & 0.19\stdv{0.04} \\
\bottomrule
\end{tabular}
\end{table}

\begin{table}[!htbp]
\centering
\caption{\textbf{CBF slack usage.} Trials (of 50) in which the QP relaxed a barrier row at least once, and the largest per-trial \texttt{slack\_max}. Configurations without a CBF (A, C) are zero by construction; on the mover-free families the CBF sees no obstacle. Source: \texttt{experiments/results\_2d/stats.json} (\texttt{slack\_usage}).}
\label{tab:slack}
\small
\begin{tabular}{lcccccc}
\toprule
Family & A & C & D & E & F & F$^{-}$ \\
\midrule
U-trap & 0 & 0 & 0 & 0 & 0 & 0 \\
Clutter & 0 & 0 & 0 & 0 & 0 & 0 \\
Dynamic & 0 & 0 & 47 (0.015) & 48 (0.015) & 48 (0.015) & 47 (0.015) \\
Narrow-dyn. & 0 & 0 & 26 (0.469) & 34 (0.469) & 34 (0.469) & 26 (0.469) \\
\bottomrule
\end{tabular}
\end{table}


\Cref{tab:full} lists every configuration and family of the benchmark, and \cref{tab:slack} the slack usage behind the discussion of \cref{sec:coord,sec:ablation}. Configurations E and F coincide exactly on U-trap and clutter because the CBF sees no obstacle in those mover-free families, so the gain schedule is irrelevant there.

\section{Notation}\label{app:notation}

\begin{table}[!htbp]
\centering
\caption{\textbf{Notation.}}
\label{tab:notation}
\small
\begin{tabular}{ll}
\toprule
Symbol & Meaning \\
\midrule
$\mathbf{x}_t=(x_t,y_t,\theta_t)$, $\mathbf{p}_t$ & robot pose, robot position \\
$\mathbf{u}_t=(v_t,\omega_t)$, $\mathcal{U}$ & control (linear, angular velocity), admissible box \\
$\sigma=(\sigma_1,\dots,\sigma_M)$ & global path from the Nav2 planner \\
$k_t,\ k^{\max}_t$ & nearest and furthest-reached path index \\
$s_t,\ W,\ e_t$ & stall counter, stall window, entrapment flag \\
$\mathcal{O}_t$; $\mathbf{p}_o, R_o, \mathbf{v}_o$ & tracked dynamic obstacles; position, radius, velocity \\
$\mathbf{P}_t,\ L,\ G(\theta)$ & look-ahead point, its offset, its Jacobian \\
$h_o,\ R^{\mathrm{eff}}_o$ & barrier for obstacle $o$, effective radius $r+R_o+m$ \\
$\alpha_t$; $\alpha_{\mathrm{base}},\alpha_{\mathrm{escape}}$ & class-$\mathcal{K}$ gain; its two scheduled values \\
$\delta,\ \rho,\ \epsilon_\delta$ & QP slack, slack weight, braking threshold \\
$\mathrm{TTC}_t,\ \tau$ & minimum time-to-collision, override threshold \\
$S^k,\ w^k,\ \lambda$ & rollout cost, softmax weight, temperature \\
$C_{\mathrm{apf}},\ C_{\mathrm{gap}}$ & escape costs (repulsion, gap attraction) \\
$\psi^\star$ & bearing of the selected free-space gap \\
\bottomrule
\end{tabular}
\end{table}

